\documentclass[12pt,a4paper]{article}

% ===========================================
% PACKAGES ET CONFIGURATION
% ===========================================

\usepackage[utf8]{inputenc}
\usepackage[T1]{fontenc}
\usepackage[french]{babel}
\usepackage{geometry}
\usepackage{graphicx}
\usepackage{xcolor}
\usepackage{tikz}
\usetikzlibrary{shapes,arrows,positioning,fit,backgrounds}
\usepackage{listings}
\usepackage{fancyhdr}
\usepackage{tocloft}
\usepackage{hyperref}
\usepackage{enumitem}
\usepackage{titlesec}
\usepackage{caption}
\usepackage{subcaption}
\usepackage{amsmath}
\usepackage{amssymb}
\usepackage{booktabs}
\usepackage{longtable}
\usepackage{array}
\usepackage{tabularx}
\usepackage{float}

% Configuration de la géométrie
\geometry{
    left=2.5cm,
    right=2.5cm,
    top=3cm,
    bottom=3cm,
    headheight=15pt
}

% Configuration des couleurs
\definecolor{darkgray}{RGB}{44,44,44}
\definecolor{lightblue}{RGB}{100,181,246}
\definecolor{lightpurple}{RGB}{186,104,200}
\definecolor{lightgreen}{RGB}{129,199,132}
\definecolor{lightorange}{RGB}{255,183,77}
\definecolor{lightpink}{RGB}{240,98,146}
\definecolor{lightcyan}{RGB}{66,165,245}
\definecolor{orange}{RGB}{255,167,38}

% Configuration des listings (code)
\lstdefinestyle{bash}{
    language=bash,
    basicstyle=\ttfamily\footnotesize,
    backgroundcolor=\color{gray!10},
    frame=single,
    breaklines=true,
    showstringspaces=false,
    commentstyle=\color{green!60!black},
    keywordstyle=\color{blue!80!black},
    stringstyle=\color{red!80!black},
    numbers=left,
    numberstyle=\tiny\color{gray},
    stepnumber=1,
    numbersep=8pt
}

\lstdefinestyle{nginx}{
    language=,
    basicstyle=\ttfamily\footnotesize,
    backgroundcolor=\color{gray!10},
    frame=single,
    breaklines=true,
    showstringspaces=false,
    commentstyle=\color{green!60!black},
    keywordstyle=\color{blue!80!black},
    stringstyle=\color{red!80!black},
    numbers=left,
    numberstyle=\tiny\color{gray},
    stepnumber=1,
    numbersep=8pt
}

\lstdefinestyle{python}{
    language=Python,
    basicstyle=\ttfamily\footnotesize,
    backgroundcolor=\color{gray!10},
    frame=single,
    breaklines=true,
    showstringspaces=false,
    commentstyle=\color{green!60!black},
    keywordstyle=\color{blue!80!black},
    stringstyle=\color{red!80!black},
    numbers=left,
    numberstyle=\tiny\color{gray},
    stepnumber=1,
    numbersep=8pt
}

\lstdefinestyle{javascript}{
    language=JavaScript,
    basicstyle=\ttfamily\footnotesize,
    backgroundcolor=\color{gray!10},
    frame=single,
    breaklines=true,
    showstringspaces=false,
    commentstyle=\color{green!60!black},
    keywordstyle=\color{blue!80!black},
    stringstyle=\color{red!80!black},
    numbers=left,
    numberstyle=\tiny\color{gray},
    stepnumber=1,
    numbersep=8pt
}

% Configuration des titres
\titleformat{\section}
  {\Large\bfseries\color{blue!80!black}}
  {\thesection}
  {1em}
  {}
  [\color{blue!80!black}\titlerule]

\titleformat{\subsection}
  {\large\bfseries\color{blue!70!black}}
  {\thesubsection}
  {1em}
  {}

\titleformat{\subsubsection}
  {\normalsize\bfseries\color{blue!60!black}}
  {\thesubsubsection}
  {1em}
  {}

% Configuration des headers/footers
\pagestyle{fancy}
\fancyhf{}
\fancyhead[L]{\leftmark}
\fancyhead[R]{Documentation Technique Doc@uthANTIC}
\fancyfoot[C]{\thepage}
\renewcommand{\headrulewidth}{0.4pt}
\renewcommand{\footrulewidth}{0pt}

% Configuration des liens
\hypersetup{
    colorlinks=true,
    linkcolor=blue!80!black,
    filecolor=magenta,
    urlcolor=cyan!80!black,
    citecolor=green!80!black,
    pdftitle={Documentation Technique : Implémentation Nginx Reverse Proxy pour Doc@uthANTIC},
    pdfauthor={Équipe de développement Doc@uthANTIC},
    pdfsubject={Configuration Nginx et Architecture Microservices},
    pdfkeywords={Nginx, Reverse Proxy, Doc@uthANTIC, Microservices, Django, Vue.js}
}

% Configuration de la table des matières
\renewcommand{\cftsecleader}{\cftdotfill{\cftdotsep}}

% ===========================================
% DÉBUT DU DOCUMENT
% ===========================================

\begin{document}

% Page de titre
\begin{titlepage}
    \centering
    \vspace*{2cm}
    
    {\Huge\bfseries Documentation Technique}\\[0.5cm]
    {\LARGE Implémentation Nginx Reverse Proxy}\\[0.3cm]
    {\LARGE pour Doc@uthANTIC}\\[2cm]
    
    \begin{tikzpicture}[scale=0.8]
        % Logo symbolique avec TikZ
        \draw[thick, blue!80!black] (0,0) rectangle (6,4);
        \node[align=center, font=\Large\bfseries] at (3,2) {Doc@uthANTIC\\Nginx\\Proxy};
        \draw[thick, orange] (-1,1.5) rectangle (0,2.5);
        \draw[thick, orange] (6,1.5) rectangle (7,2.5);
        \node[font=\small] at (-0.5,2) {Client};
        \node[font=\small] at (6.5,2) {Services};
        \draw[->, thick, orange] (0,2) -- (1,2);
        \draw[->, thick, orange] (5,2) -- (6,2);
    \end{tikzpicture}\\[2cm]
    
    {\large De l'installation à la production}\\[1cm]
    {\large Résolution des problèmes CORS, CSRF et permissions}\\[3cm]
    
    {\large\textbf{Équipe de développement Doc@uthANTIC}}\\[0.5cm]
    {\large\today}
\end{titlepage}

% Table des matières
\tableofcontents
\newpage

% ===========================================
% SECTION 1: VUE D'ENSEMBLE
% ===========================================

\section{Vue d'ensemble du projet}

\subsection{Objectif principal}

Transformer l'architecture Doc@uthANTIC d'un système exposant directement les ports des services vers une architecture sécurisée utilisant un reverse proxy Nginx, éliminant l'exposition des ports explicites dans les URLs.

\subsection{Problématique initiale}

\begin{itemize}[leftmargin=*]
    \item \textbf{URLs exposant les ports :} \texttt{https://192.168.4.131:8000}, \texttt{https://192.168.4.131:8001}, etc.
    \item \textbf{Services accessibles directement :} Risques de sécurité
    \item \textbf{Configuration CORS complexe :} Gestion manuelle pour chaque service
    \item \textbf{Maintenance difficile :} Changement d'URL implique modification du code
\end{itemize}

\subsection{Solution implémentée}

\begin{itemize}[leftmargin=*]
    \item \textbf{Reverse proxy Nginx :} Point d'entrée unique sur les ports 80/443
    \item \textbf{Services en localhost :} Sécurisation par isolation réseau
    \item \textbf{URLs simplifiées :} \texttt{https://192.168.4.131/api/}, \texttt{https://192.168.4.131/gateway/}, etc.
    \item \textbf{SSL centralisé :} Gestion des certificats uniquement sur Nginx
\end{itemize}

% ===========================================
% SECTION 2: ARCHITECTURE INITIALE
% ===========================================

\section{Architecture initiale}

\subsection{Diagramme de l'architecture avant Nginx}

\begin{figure}[H]
    \centering
    \begin{tikzpicture}[
        node distance=1.5cm and 2cm,
        box/.style={rectangle, draw, fill=darkgray, text=white, minimum width=2.5cm, minimum height=1cm, font=\small\bfseries},
        client/.style={rectangle, draw, fill=blue!20, minimum width=2cm, minimum height=0.8cm, font=\small},
        arrow/.style={->, thick}
    ]
        % Client
        \node[client] (client) {Client Web Browser};
        
        % Services avant Nginx
        \node[box, below=2cm of client, xshift=-4cm] (frontend) {Vue.js Frontend\\:8080};
        \node[box, right=of frontend] (django) {Django Backend\\:8000};
        \node[box, right=of django] (gateway) {API Gateway\\:8001};
        \node[box, below=of frontend] (certms) {Certificat MS\\:8002};
        \node[box, below=of django] (signms) {Signature MS\\:8003};
        
        % Connexions directes (problématiques)
        \draw[arrow, red] (client) to[out=-135,in=90] node[left, font=\tiny] {:8080} (frontend);
        \draw[arrow, red] (client) to[out=-90,in=90] node[above, font=\tiny] {:8000} (django);
        \draw[arrow, red] (client) to[out=-45,in=90] node[right, font=\tiny] {:8001} (gateway);
        
        % Connexions inter-services
        \draw[arrow] (frontend) -- node[above, font=\tiny] {:8000} (django);
        \draw[arrow] (frontend) to[out=0,in=180] node[above, font=\tiny] {:8001} (gateway);
        \draw[arrow] (gateway) -- node[left, font=\tiny] {:8002} (certms);
        \draw[arrow] (gateway) -- node[right, font=\tiny] {:8003} (signms);
        
        % Titre
        \node[above=0.5cm of client, font=\Large\bfseries] {Services Doc@uthANTIC (AVANT)};
        
        % Légende
        \node[below=3cm of signms, font=\small, text=red] {❌ Exposition directe des services - Risques de sécurité};
    \end{tikzpicture}
    \caption{Architecture avant l'implémentation du reverse proxy Nginx}
    \label{fig:architecture_avant}
\end{figure}

\subsection{Services et ports initiaux}

\begin{table}[H]
\centering
\caption{Configuration des services avant Nginx}
\label{tab:services_initiaux}
\begin{tabularx}{\textwidth}{|l|c|X|X|}
\hline
\textbf{Service} & \textbf{Port} & \textbf{URL d'accès} & \textbf{Fonction} \\
\hline
Frontend Vue.js & 8080 & \texttt{https://192.168.4.131:8080} & Interface utilisateur \\
\hline
Backend Django & 8000 & \texttt{https://192.168.4.131:8000} & API REST, Admin \\
\hline
API Gateway & 8001 & \texttt{https://192.168.4.131:8001} & Routage vers microservices \\
\hline
Microservice Certificat & 8002 & \texttt{https://192.168.4.131:8002} & Traitement certificats \\
\hline
Microservice Signature & 8003 & \texttt{https://192.168.4.131:8003} & Signature documents \\
\hline
\end{tabularx}
\end{table}

\subsection{Problèmes identifiés}

\subsubsection{🔴 Sécurité}
\begin{itemize}[leftmargin=*]
    \item \textbf{Exposition directe des services :} Chaque service accessible depuis l'extérieur
    \item \textbf{Gestion SSL multiple :} Certificats sur chaque service
    \item \textbf{Attaques directes possibles :} Contournement du frontend
\end{itemize}

\subsubsection{🔴 Maintenance}
\begin{itemize}[leftmargin=*]
    \item \textbf{URLs hardcodées :} Ports dans le code frontend
    \item \textbf{Configuration CORS répétée :} Sur chaque service
    \item \textbf{Déploiement complexe :} Gestion de 5 services distincts
\end{itemize}

\subsubsection{🔴 Évolutivité}
\begin{itemize}[leftmargin=*]
    \item \textbf{Ajout de services difficile :} Nouveau port à exposer
    \item \textbf{Load balancing impossible :} Pas de point central
    \item \textbf{Monitoring dispersé :} Logs sur chaque service
\end{itemize}

% ===========================================
% SECTION 3: INSTALLATION NGINX
% ===========================================

\section{Installation et configuration Nginx}

\subsection{Installation sur Ubuntu}

La première étape consistait à installer Nginx sur le système Ubuntu :

\begin{lstlisting}[style=bash, caption={Installation de Nginx sur Ubuntu}]
# Mise à jour des paquets
sudo apt update

# Installation de Nginx
sudo apt install nginx -y

# Vérification de l'installation
nginx -v
# nginx version: nginx/1.18.0 (Ubuntu)

# Démarrage et activation automatique
sudo systemctl start nginx
sudo systemctl enable nginx

# Vérification du statut
sudo systemctl status nginx
\end{lstlisting}

\subsection{Création des certificats SSL}

Pour sécuriser les communications, nous avons réutilisé les certificats existants du projet :

\begin{lstlisting}[style=bash, caption={Configuration des certificats SSL}]
# Création du répertoire SSL pour Nginx
sudo mkdir -p /etc/nginx/ssl/docuthanatic/

# Copie des certificats existants
sudo cp backend/django-project/ssl/cert.pem /etc/nginx/ssl/docuthanatic/
sudo cp backend/django-project/ssl/key.pem /etc/nginx/ssl/docuthanatic/

# Vérification des permissions
sudo chown root:root /etc/nginx/ssl/docuthanatic/*
sudo chmod 600 /etc/nginx/ssl/docuthanatic/key.pem
sudo chmod 644 /etc/nginx/ssl/docuthanatic/cert.pem
\end{lstlisting}

\subsection{Configuration initiale Nginx}

Création du fichier de configuration principal \texttt{/etc/nginx/sites-available/docuthanatic} :

\begin{lstlisting}[style=nginx, caption={Configuration Nginx pour Doc@uthANTIC}, label={lst:nginx_config}]
# Configuration Nginx pour Doc@uthANTIC
# Ce fichier dit à Nginx comment rediriger les requêtes vers nos services

server {
    # Écouter sur le port 80 (HTTP) et 443 (HTTPS)
    listen 80;
    listen 443 ssl;
    
    # Nom du serveur (votre adresse IP)
    server_name 192.168.4.131;
    
    # Configuration SSL (certificats pour HTTPS)
    ssl_certificate /etc/nginx/ssl/docuthanatic/cert.pem;
    ssl_certificate_key /etc/nginx/ssl/docuthanatic/key.pem;
    
    # Redirection automatique HTTP vers HTTPS
    if ($scheme = http) {
        return 301 https://$server_name$request_uri;
    }
    
    # Configuration SSL sécurisée
    ssl_protocols TLSv1.2 TLSv1.3;
    ssl_ciphers HIGH:!aNULL:!MD5;
    ssl_prefer_server_ciphers on;
    
    # Headers de sécurité
    add_header X-Frame-Options DENY;
    add_header X-Content-Type-Options nosniff;
    add_header X-XSS-Protection "1; mode=block";
    
    # Gestion des gros fichiers (pour les PDFs)
    client_max_body_size 50M;
    
    # ROUTE PRINCIPALE : Frontend Vue.js
    location / {
        proxy_pass https://127.0.0.1:8080;
        proxy_set_header Host $host;
        proxy_set_header X-Real-IP $remote_addr;
        proxy_set_header X-Forwarded-For $proxy_add_x_forwarded_for;
        proxy_set_header X-Forwarded-Proto $scheme;
        
        # Support WebSocket (si besoin)
        proxy_http_version 1.1;
        proxy_set_header Upgrade $http_upgrade;
        proxy_set_header Connection "upgrade";
    }
    
    # ROUTE API : Django Backend
    location /api/ {
        proxy_pass https://127.0.0.1:8000;
        proxy_set_header Host $host;
        proxy_set_header X-Real-IP $remote_addr;
        proxy_set_header X-Forwarded-For $proxy_add_x_forwarded_for;
        proxy_set_header X-Forwarded-Proto $scheme;
    }
    
    # ROUTE ADMIN : Django Admin
    location /admin/ {
        proxy_pass https://127.0.0.1:8000;
        proxy_set_header Host $host;
        proxy_set_header X-Real-IP $remote_addr;
        proxy_set_header X-Forwarded-For $proxy_add_x_forwarded_for;
        proxy_set_header X-Forwarded-Proto $scheme;
    }
    
    # ROUTE GATEWAY : API Gateway FastAPI
    location /gateway/ {
        rewrite ^/gateway/(.*)$ /$1 break;
        proxy_pass https://127.0.0.1:8001;
        proxy_set_header Host $host;
        proxy_set_header X-Real-IP $remote_addr;
        proxy_set_header X-Forwarded-For $proxy_add_x_forwarded_for;
        proxy_set_header X-Forwarded-Proto $scheme;
    }
    
    # ROUTE CERTIFICATS : Microservice Certificats
    location /cert/ {
        rewrite ^/cert/(.*)$ /$1 break;
        proxy_pass https://127.0.0.1:8002;
        proxy_set_header Host $host;
        proxy_set_header X-Real-IP $remote_addr;
        proxy_set_header X-Forwarded-For $proxy_add_x_forwarded_for;
        proxy_set_header X-Forwarded-Proto $scheme;
    }
    
    # ROUTE SIGNATURE : Microservice Signature
    location /sign/ {
        rewrite ^/sign/(.*)$ /$1 break;
        proxy_pass https://127.0.0.1:8003;
        proxy_set_header Host $host;
        proxy_set_header X-Real-IP $remote_addr;
        proxy_set_header X-Forwarded-For $proxy_add_x_forwarded_for;
        proxy_set_header X-Forwarded-Proto $scheme;
    }
    
    # ROUTE MEDIA : Fichiers uploadés (PDFs, images)
    location /media/ {
        proxy_pass https://127.0.0.1:8000;
        proxy_set_header Host $host;
        proxy_set_header X-Real-IP $remote_addr;
        proxy_set_header X-Forwarded-For $proxy_add_x_forwarded_for;
        proxy_set_header X-Forwarded-Proto $scheme;
    }
    
    # ROUTE STATIC : Fichiers statiques (CSS, JS, images)
    location /static/ {
        proxy_pass https://127.0.0.1:8000;
        proxy_set_header Host $host;
        proxy_set_header X-Real-IP $remote_addr;
        proxy_set_header X-Forwarded-For $proxy_add_x_forwarded_for;
        proxy_set_header X-Forwarded-Proto $scheme;
    }
    
    # Logs pour debug
    access_log /var/log/nginx/docuthanatic_access.log;
    error_log /var/log/nginx/docuthanatic_error.log;
}
\end{lstlisting}

\subsection{Activation de la configuration}

\begin{lstlisting}[style=bash, caption={Activation du site Nginx}]
# Création du lien symbolique pour activer le site
sudo ln -s /etc/nginx/sites-available/docuthanatic /etc/nginx/sites-enabled/

# Désactivation du site par défaut
sudo rm /etc/nginx/sites-enabled/default

# Test de la configuration
sudo nginx -t

# Rechargement de Nginx
sudo systemctl reload nginx
\end{lstlisting}

\subsection{Routage mis en place}

Le nouveau routage Nginx transforme les URLs comme suit :

\begin{table}[H]
\centering
\caption{Transformation des URLs par Nginx}
\label{tab:routage_nginx}
\begin{tabularx}{\textwidth}{|X|X|X|}
\hline
\textbf{URL d'entrée} & \textbf{Service de destination} & \textbf{Fonction} \\
\hline
\texttt{https://192.168.4.131/} & \texttt{https://127.0.0.1:8080} & Frontend Vue.js \\
\hline
\texttt{https://192.168.4.131/api/} & \texttt{https://127.0.0.1:8000/api/} & API Django \\
\hline
\texttt{https://192.168.4.131/admin/} & \texttt{https://127.0.0.1:8000/admin/} & Django Admin \\
\hline
\texttt{https://192.168.4.131/gateway/} & \texttt{https://127.0.0.1:8001/} & API Gateway \\
\hline
\texttt{https://192.168.4.131/cert/} & \texttt{https://127.0.0.1:8002/} & Microservice Certificats \\
\hline
\texttt{https://192.168.4.131/sign/} & \texttt{https://127.0.0.1:8003/} & Microservice Signature \\
\hline
\texttt{https://192.168.4.131/media/} & \texttt{https://127.0.0.1:8000/media/} & Fichiers média \\
\hline
\texttt{https://192.168.4.131/static/} & \texttt{https://127.0.0.1:8000/static/} & Fichiers statiques \\
\hline
\end{tabularx}
\end{table}

% ===========================================
% SECTION 4: CONFIGURATION SERVICES LOCALHOST
% ===========================================

\section{Configuration des services pour localhost}

\subsection{Objectif de sécurisation}

Une fois Nginx configuré, il fallait modifier tous les services backend pour qu'ils n'écoutent plus que sur localhost (\texttt{127.0.0.1}) au lieu de l'adresse IP publique, garantissant ainsi qu'ils ne soient accessibles que via le reverse proxy.

\subsection{Diagramme de l'architecture réseau modifiée}

\begin{figure}[H]
    \centering
    \begin{tikzpicture}[
        node distance=1.5cm and 2cm,
        nginx/.style={rectangle, draw, fill=orange, minimum width=2.5cm, minimum height=1cm, font=\small\bfseries},
        box/.style={rectangle, draw, fill=darkgray, text=white, minimum width=2.5cm, minimum height=1cm, font=\small\bfseries},
        client/.style={rectangle, draw, fill=blue!20, minimum width=2cm, minimum height=0.8cm, font=\small},
        arrow/.style={->, thick}
    ]
        % Client
        \node[client] (client) {Client Web Browser};
        
        % Nginx (nouveau)
        \node[nginx, below=1.5cm of client] (nginx) {Nginx Reverse Proxy\\:80/:443};
        
        % Services après Nginx (localhost)
        \node[box, below=2cm of nginx, xshift=-4cm] (frontend) {Vue.js Frontend\\127.0.0.1:8080};
        \node[box, right=of frontend] (django) {Django Backend\\127.0.0.1:8000};
        \node[box, right=of django] (gateway) {API Gateway\\127.0.0.1:8001};
        \node[box, below=of frontend] (certms) {Certificat MS\\127.0.0.1:8002};
        \node[box, below=of django] (signms) {Signature MS\\127.0.0.1:8003};
        
        % Connexions via Nginx (sécurisées)
        \draw[arrow, green!80!black] (client) -- node[right, font=\tiny] {:443} (nginx);
        \draw[arrow, blue] (nginx) to[out=-135,in=90] node[left, font=\tiny] {:8080} (frontend);
        \draw[arrow, blue] (nginx) to[out=-90,in=90] node[above, font=\tiny] {:8000} (django);
        \draw[arrow, blue] (nginx) to[out=-45,in=90] node[right, font=\tiny] {:8001} (gateway);
        
        % Connexions inter-services (localhost)
        \draw[arrow] (frontend) -- node[above, font=\tiny] {:8000} (django);
        \draw[arrow] (frontend) to[out=0,in=180] node[above, font=\tiny] {:8001} (gateway);
        \draw[arrow] (gateway) -- node[left, font=\tiny] {:8002} (certms);
        \draw[arrow] (gateway) -- node[right, font=\tiny] {:8003} (signms);
        
        % Titre
        \node[above=0.5cm of client, font=\Large\bfseries] {Services Doc@uthANTIC (APRÈS)};
        
        % Légende
        \node[below=3cm of signms, font=\small, text=green!80!black] {✅ Point d'entrée unique - Sécurité renforcée};
    \end{tikzpicture}
    \caption{Architecture après implémentation du reverse proxy Nginx}
    \label{fig:architecture_apres}
\end{figure}

\subsection{Modifications des scripts de démarrage}

Les scripts de démarrage des services ont été modifiés pour utiliser \texttt{127.0.0.1} :

\begin{lstlisting}[style=bash, caption={Modifications des scripts de démarrage}]
# Django Backend - AVANT
python3 manage.py runserver_plus --cert-file ssl/cert.pem \
    --key-file ssl/key.pem 192.168.4.131:8000 --insecure

# Django Backend - APRÈS  
python3 manage.py runserver_plus --cert-file ssl/cert.pem \
    --key-file ssl/key.pem 127.0.0.1:8000 --insecure

# API Gateway FastAPI - AVANT
uvicorn main:app --host 192.168.4.131 --port 8001 \
    --ssl-keyfile ssl/key.pem --ssl-certfile ssl/cert.pem

# API Gateway FastAPI - APRÈS
uvicorn main:app --host 127.0.0.1 --port 8001 \
    --ssl-keyfile ssl/key.pem --ssl-certfile ssl/cert.pem

# Frontend Vue.js - AVANT
npm run serve -- --host 192.168.4.131 --port 8080 --https

# Frontend Vue.js - APRÈS
npm run serve -- --host 127.0.0.1 --port 8080 --https
\end{lstlisting}

% ===========================================
% SECTION 5: RÉSOLUTION PROBLÈMES CORS
% ===========================================

\section{Résolution des problèmes CORS}

\subsection{Problème CORS identifié}

Après la configuration initiale d'Nginx, les requêtes depuis le frontend Vue.js vers l'API Django étaient bloquées par les politiques CORS. Le navigateur affichait cette erreur :

\begin{lstlisting}[style=bash, caption={Erreur CORS dans la console du navigateur}]
Access to XMLHttpRequest at 'https://192.168.4.131/api/users/auth/with-organization/' 
from origin 'https://192.168.4.131' has been blocked by CORS policy: 
No 'Access-Control-Allow-Origin' header is present on the requested resource.
\end{lstlisting}

\subsection{Analyse du problème}

\begin{itemize}[leftmargin=*]
    \item \textbf{Origine des requêtes :} Frontend sur \texttt{https://192.168.4.131}
    \item \textbf{Destination :} API Django via \texttt{https://192.168.4.131/api/}
    \item \textbf{Problème :} Django configuré pour \texttt{127.0.0.1:8000}, pas \texttt{192.168.4.131}
    \item \textbf{Solution :} Ajouter \texttt{https://192.168.4.131} aux origins CORS autorisées
\end{itemize}

\subsection{Configuration Django CORS}

\subsubsection{Modification du fichier settings.py}

\begin{lstlisting}[style=python, caption={Configuration CORS Django}, label={lst:django_cors}]
# backend/django-project/certisign_project/settings.py

# Configuration CORS mise à jour pour Nginx
CORS_ALLOWED_ORIGINS = [
    "https://127.0.0.1:8080",  # Frontend direct (développement)
    "https://192.168.4.131",   # Nginx reverse proxy (production)
    "https://192.168.4.131:8080",  # Frontend avec port (fallback)
]

# Autoriser les cookies et headers d'authentification
CORS_ALLOW_CREDENTIALS = True

# Headers autorisés pour CORS
CORS_ALLOW_HEADERS = [
    'accept',
    'accept-encoding',
    'authorization',
    'content-type',
    'dnt',
    'origin',
    'user-agent',
    'x-csrftoken',
    'x-requested-with',
    'x-forwarded-for',
    'x-forwarded-proto',
    'x-real-ip',
]

# Méthodes HTTP autorisées
CORS_ALLOW_METHODS = [
    'DELETE',
    'GET',
    'OPTIONS',
    'PATCH',
    'POST',
    'PUT',
]
\end{lstlisting}

\subsection{Tests de validation CORS}

\begin{lstlisting}[style=bash, caption={Tests de connectivité CORS}]
# Test d'accès à l'API via Nginx
curl -k -H "Origin: https://192.168.4.131" \
     -H "Access-Control-Request-Method: GET" \
     -H "Access-Control-Request-Headers: authorization" \
     -X OPTIONS \
     https://192.168.4.131/api/users/organizations/

# Réponse attendue avec headers CORS
# Access-Control-Allow-Origin: https://192.168.4.131
# Access-Control-Allow-Methods: GET, POST, PUT, PATCH, DELETE, OPTIONS
# Access-Control-Allow-Headers: authorization, content-type, ...
\end{lstlisting}

% ===========================================
% SECTION 6: RÉSOLUTION PROBLÈMES CSRF
% ===========================================

\section{Résolution des problèmes CSRF}

\subsection{Problème CSRF identifié}

Après la résolution des problèmes CORS, un nouveau problème est apparu : les requêtes d'authentification échouaient avec une erreur CSRF :

\begin{lstlisting}[style=bash, caption={Erreur CSRF dans l'API Django}]
POST /api/users/auth/with-organization/ HTTP/1.1" 403 58
{"detail": "CSRF Failed: CSRF token missing or incorrect."}
\end{lstlisting}

\subsection{Analyse du problème CSRF}

Le problème venait du fait que les tokens CSRF générés par Django n'étaient pas correctement transmis à travers le reverse proxy Nginx.

\subsubsection{Origine du problème}
\begin{itemize}[leftmargin=*]
    \item \textbf{Token CSRF généré par Django :} Sur \texttt{127.0.0.1:8000}
    \item \textbf{Requêtes envoyées via Nginx :} Sur \texttt{192.168.4.131}
    \item \textbf{Headers CSRF manquants :} Nginx ne transmettait pas les tokens
    \item \textbf{Cookies de domaine :} Incompatibilité entre les domaines
\end{itemize}

\subsection{Solution 1 : Configuration Nginx pour CSRF}

\subsubsection{Modification de : \texttt{/etc/nginx/sites-available/docuthanatic}}

\textbf{Section API Django mise à jour :}

\begin{lstlisting}[style=nginx, caption={Configuration Nginx pour transmission CSRF}]
# ROUTE API : Django Backend avec support CSRF
location /api/ {
    proxy_pass https://127.0.0.1:8000;
    proxy_set_header Host $host;
    proxy_set_header X-Real-IP $remote_addr;
    proxy_set_header X-Forwarded-For $proxy_add_x_forwarded_for;
    proxy_set_header X-Forwarded-Proto $scheme;
    
    # IMPORTANT : Transmission des tokens CSRF
    proxy_set_header X-CSRFToken $cookie_csrftoken;
    proxy_set_header Referer $scheme://$host$request_uri;
    
    # Support des cookies de session et CSRF
    proxy_cookie_domain 127.0.0.1 $host;
    proxy_cookie_path / /;
}
\end{lstlisting}

\subsection{Solution 2 : Exemption CSRF pour l'endpoint JWT}

\subsubsection{Modification de views\_auth.py}

Puisque l'endpoint d'authentification utilise JWT (qui est stateless), nous avons ajouté l'exemption CSRF :

\begin{lstlisting}[style=python, caption={Exemption CSRF pour l'authentification JWT}]
# backend/django-project/users/views_auth.py

from django.views.decorators.csrf import csrf_exempt
from django.utils.decorators import method_decorator

@method_decorator(csrf_exempt, name='dispatch')
class AuthWithOrganizationView(APIView):
    """
    Authentification avec organisation spécifique
    Exempt de CSRF car utilise JWT (stateless)
    """
    permission_classes = [AllowAny]
    
    def post(self, request):
        email = request.data.get('email')
        password = request.data.get('password')
        organization_id = request.data.get('organization_id')
        
        # Logique d'authentification...
        # ...reste du code inchangé...
\end{lstlisting}

\subsection{Tests de validation CSRF}

\begin{lstlisting}[style=bash, caption={Tests d'authentification après résolution CSRF}]
# Test d'authentification via Nginx
curl -k -X POST \
     -H "Content-Type: application/json" \
     -H "Origin: https://192.168.4.131" \
     -d '{
       "email": "test@example.com",
       "password": "testpass",
       "organization_id": 1
     }' \
     https://192.168.4.131/api/users/auth/with-organization/

# Réponse attendue avec JWT token
# {
#   "access": "eyJ0eXAiOiJKV1QiLCJhbGciOiJIUzI1NiJ9...",
#   "refresh": "eyJ0eXAiOiJKV1QiLCJhbGciOiJIUzI1NiJ9...",
#   "user": {...},
#   "organization": {...}
# }
\end{lstlisting}
\end{document} 